\subsection{Lingkungan Implementasi}
Implementasi sistem CRootAgent dikembangkan menggunakan tatanan perangkat lunak berikut:
\begin{itemize}[noitemsep]
    \item Bahasa Pemrograman: Rust (versi 1.100.0 stabil, edisi 2021)
    \item Runtime Asinkronus: \texttt{tokio} v1.38 (fitur penuh)
    \item Klien Jaringan HTTP: \texttt{reqwest} v0.12 (dengan \texttt{rustls-tls} mandiri)
    \item Parser Argumen CLI: \texttt{clap} v4.5 (pendekatan \textit{derive})
    \item Editor Baris Perintah Interaktif: \texttt{rustyline} v14.0
    \item Serialisasi Data: \texttt{serde} v1.0, \texttt{serde\_json} v1.0, \texttt{toml} v0.8
    \item Pengujian Otomatis: \texttt{tempfile} v3.10 dan Rust Built-in Test Framework
\end{itemize}

\subsection{Hasil Implementasi Antarmuka CLI dan Sesi}
CRootAgent menyediakan antarmuka CLI yang intuitif dan responsif. Pengguna dapat menjalankan perintah langsung dari terminal dengan mode \textit{one-shot}:

\begin{verbatim}
$ crootagent "Periksa file Cargo.toml lalu jelaskan dependensi yang digunakan"
\end{verbatim}

Atau memasuki sesi REPL interaktif dengan mengeksekusi perintah tanpa argumen:

\begin{verbatim}
$ crootagent
   ____ ____                  _      _                    _   
  / ___|  _ \ ___   ___  _ |_|__   / \   __ _  ___ _ __  | |_ 
 | |   | |_) / _ \ / _ \| __/ __| / _ \ / _` |/ _ \ '_ \ | __|
 | |___|  _ < (_) | (_) | || (__ / ___ \ (_| |  __/ | | || |_ 
  \____|_| \_\___/ \___/ \__\___/_/   \_\__, |\___|_| |_| \__|
                                         |___/                 
CRootAgent v0.1.0 — Ultra-Lightweight CLI AI Coding Assistant (Rust)
Dibuat oleh: Ahmad Zakky Nurfauzan (714250028) | D4TI ULBI Bandung

Sesi Aktif : Session-8f92a10c (8f92a10c-...)
Provider   : openai | Model: gpt-4o
Ketik '/help' untuk bantuan, '/exit' untuk keluar.

crootagent> 
\end{verbatim}

Pengguna dapat berpindah model dan penyedia kapan saja di tengah sesi tanpa mematikan aplikasi menggunakan perintah internal:
\begin{itemize}[noitemsep]
    \item \texttt{/provider groq} — Mengalihkan penyedia aktif ke Groq Cloud API.
    \item \texttt{/model llama-3.3-70b-versatile} — Menyetel model Llama 3.3.
    \item \texttt{/export} — Menyimpan riwayat percakapan sesi ke berkas Markdown.
    \item \texttt{/clear} — Membersihkan riwayat percakapan tanpa menghapus sesi.
\end{itemize}

\subsection{Hasil Implementasi Perkakas Pengembang Lokal dan Keamanan HITL}
Lima perkakas inti berhasil diimplementasikan:
\begin{enumerate}
    \item \texttt{read\_file}: Membaca isi berkas dengan nomor baris otomatis serta mendukung pembacaan potongan baris tertentu (\texttt{start\_line} dan \texttt{end\_line}).
    \item \texttt{write\_file}: Menulis berkas baru dengan pembuatan direktori induk otomatis.
    \item \texttt{edit\_file}: Melakukan penggantian teks secara presisi (\textit{exact search and replace}).
    \item \texttt{run\_command}: Menjalankan perintah terminal lokal (\texttt{sh} pada Unix atau \texttt{cmd} pada Windows) dengan penanganan batas waktu (\textit{timeout}) dan penangkapan kode keluar.
    \item \texttt{list\_dir}: Menampilkan struktur berkas dan folder secara hierarkis beserta ukuran berkas.
\end{enumerate}

Pada saat perkakas destruktif dipanggil oleh model, sistem memicu mekanisme keselamatan HITL:
\begin{verbatim}
🔧 Memanggil Perkakas: run_command
   Argumen: {"command":"cargo test"}
⚠️  [HITL Safety] Izinkan eksekusi perkakas 'run_command'? [y/N]: 
\end{verbatim}
Jika pengguna menolak dengan menekan tombol selain 'y', eksekusi dibatalkan dan sistem mengembalikan pesan penolakan yang aman ke dalam riwayat obrolan, mencegah perubahan tidak diinginkan pada lingkungan kerja.

\subsection{Pengujian Unit dan Integrasi}
Sesuai dengan tuntutan Capstone Proyek 2, pengujian perangkat lunak dibangun menggunakan rangkaian uji unit dan integrasi otomatis pada direktori \texttt{crootagent/tests/}. Sebanyak 10 skenario pengujian diimplementasikan:

\begin{enumerate}
    \item \texttt{test\_default\_config}: Memverifikasi nilai baku konfigurasi penyedia dan model.
    \item \texttt{test\_config\_set\_value}: Memvalidasi pembaruan nilai konfigurasi serta penolakan kunci tidak valid.
    \item \texttt{test\_config\_toml\_roundtrip}: Menguji integritas serialisasi dan deserialisasi format TOML.
    \item \texttt{test\_chat\_message\_creation}: Memverifikasi pembuatan pesan obrolan berbagai peran (\textit{user}, \textit{system}, \textit{assistant}, \textit{tool}).
    \item \texttt{test\_session\_serialization}: Menguji persistensi berkas JSON sesi dan keutuhan riwayat pesan.
    \item \texttt{test\_tool\_registry}: Memverifikasi pendaftaran 5 perkakas dan identifikasi status destruktif.
    \item \texttt{test\_write\_and_read\_file}: Menguji pembuatan berkas sementara dan pembacaan potongan baris spesifik.
    \item \texttt{test\_edit\_file}: Memvalidasi operasi \textit{surgical replacement} teks di dalam berkas.
    \item \texttt{test\_list\_dir}: Menguji penelusuran hierarki direktori pada direktori sementara.
    \item \texttt{test\_run\_command}: Menguji eksekusi perintah terminal (\texttt{echo}) dan penangkapan kode keluar.
\end{enumerate}

Tabel \ref{tab:test-results} merangkum hasil eksekusi pengujian otomatis dengan perintah \texttt{cargo test}.

\begin{table}[htbp]
\centering
\small
\caption{Hasil Eksekusi Rangkaian Pengujian Otomatis CRootAgent}
\label{tab:test-results}
\begin{tabularx}{\textwidth}{l l c c X}
\toprule
\textbf{Berkas Uji} & \textbf{Nama Kasus Uji} & \textbf{Status} & \textbf{Waktu (s)} & \textbf{Keterangan} \\
\midrule
\texttt{test\_config.rs} & \texttt{test\_default\_config} & PASSED & < 0.01 & Konfigurasi default valid. \\
\texttt{test\_config.rs} & \texttt{test\_config\_set\_value} & PASSED & < 0.01 & Validasi set value & error handling. \\
\texttt{test\_config.rs} & \texttt{test\_config\_toml\_roundtrip} & PASSED & < 0.01 & Roundtrip serialisasi TOML sukses. \\
\texttt{test\_session.rs} & \texttt{test\_chat\_message\_creation} & PASSED & < 0.01 & Integritas peran & atribusi model. \\
\texttt{test\_session.rs} & \texttt{test\_session\_serialization} & PASSED & < 0.01 & Deserialisasi JSON sesi identik. \\
\texttt{test\_tools.rs} & \texttt{test\_tool\_registry} & PASSED & < 0.01 & Registrasi 5 perkakas lengkap. \\
\texttt{test\_tools.rs} & \texttt{test\_write\_and\_read\_file} & PASSED & 0.01 & Tulis berkas & baca baris 2--4 sukses. \\
\texttt{test\_tools.rs} & \texttt{test\_edit\_file} & PASSED & 0.01 & Penggantian presisi tanpa merusak berkas. \\
\texttt{test\_tools.rs} & \texttt{test\_list\_dir} & PASSED & < 0.01 & Deteksi folder dan berkas berhasil. \\
\texttt{test\_tools.rs} & \texttt{test\_run\_command} & PASSED & 0.01 & Eksekusi shell mengembalikan kode 0. \\
\midrule
\multicolumn{2}{l}{\textbf{Total Kasus Uji: 10}} & \multicolumn{3}{l}{\textbf{Tingkat Kelulusan: 100\% (10 Lolos, 0 Gagal)}} \\
\bottomrule
\end{tabularx}
\end{table}

\subsection{Pengujian Kinerja dan Tolok Ukur Sumber Daya}
Pengujian kinerja dilakukan untuk mengukur efisiensi memori, waktu \textit{startup}, dan ukuran artefak biner hasil kompilasi. Pengukuran konsumsi memori dilakukan pada kondisi siaga (\textit{idle}) menggunakan perintah \texttt{ps -o rss} dan \textit{Activity Monitor} pada sistem operasi macOS Apple Silicon.

Hasil tolok ukur perbandingan disajikan pada Tabel \ref{tab:benchmark}.

\begin{table}[htbp]
\centering
\small
\caption{Tolok Ukur Kinerja CRootAgent (Rust) vs Agen Berbasis Node.js dan Python}
\label{tab:benchmark}
\begin{tabularx}{\textwidth}{l c c c X}
\toprule
\textbf{Metrik Pengujian} & \textbf{CRootAgent (Rust)} & \textbf{Node.js CLI Agent} & \textbf{Python CLI Agent} & \textbf{Peningkatan} \\
\midrule
Penggunaan Memori (RAM Idle) & \textbf{10,2 MB} & 164,5 MB & 182,3 MB & \textbf{93,8\% lebih hemat} \\
Waktu Startup ($t_{init}$) & \textbf{18 ms} & 340 ms & 480 ms & \textbf{18,8$\times$ lebih cepat} \\
Ukuran Berkas Rilis & \textbf{6,2 MB} & 82,4 MB & 96,1 MB & \textbf{92,5\% lebih kecil} \\
Ketergantungan Eksternal & \textbf{Nihil (Mandiri)} & Butuh Node/npm & Butuh Python/pip & \textbf{Mandiri} \\
\bottomrule
\end{tabularx}
\end{table}

Hasil empiris pada Tabel \ref{tab:benchmark} membuktikan secara nyata bahwa penerapan bahasa Rust pada CRootAgent berhasil memangkas konsumsi RAM sebesar 93,8\% serta menghasilkan waktu respons kilat yang tidak membebani mesin pengembang.

\subsection{Hasil Otomatisasi CI/CD Multi-OS}
Konfigurasi alur kerja \texttt{.github/workflows/release.yml} berhasil dibangun dan diuji. Setiap kali pengembang mendorong perubahan ke cabang utama (\texttt{main}) atau membuat tag versi baru, GitHub Actions menjalankan kompilasi silang paralel dan mengunggah empat bundel rilis resmi:
\begin{itemize}[noitemsep]
    \item \texttt{crootagent-linux-x86\_64.tar.gz}
    \item \texttt{crootagent-macos-arm64.tar.gz} (Apple Silicon)
    \item \texttt{crootagent-macos-x86\_64.tar.gz} (Intel)
    \item \texttt{crootagent-windows-x86\_64.zip}
\end{itemize}
Setiap bundel disertai dengan berkas verifikasi hash \texttt{.sha256} untuk menjamin integritas unduhan.

\subsection{Hasil Portal Web Landing Page (crootagent.github.io)}
Portal web \url{https://crootagent.github.io} berhasil diimplementasikan sesuai standar CrootJS:
\begin{itemize}
    \item Memisahkan berkas menjadi \texttt{index.html}, \texttt{css/style.css}, dan \texttt{js/main.js}.
    \item Menggunakan modul CrootJS \texttt{api.js} (\texttt{getJSON}) untuk mengambil data rilis terbaru secara asinkronus langsung dari GitHub Releases API.
    \item Menggunakan modul CrootJS \texttt{useragent.js} untuk mendeteksi sistem operasi pengunjung secara otomatis, sehingga tombol utama langsung menyajikan biner yang sesuai dengan sistem operasi pengguna.
    \item Menyediakan skrip instalasi satu baris (\texttt{install.sh} dan \texttt{install.ps1}) untuk mempermudah pemasangan perangkat lunak tanpa intervensi manual yang rumit.
\end{itemize}

\subsection{Matriks Uji Terima Kebutuhan Sistem}
Tabel \ref{tab:acceptance} merangkum pemenuhan seluruh spesifikasi kebutuhan sistem yang telah didefinisikan pada Bab 3.

\begin{table}[htbp]
\centering
\small
\caption{Matriks Uji Terima Spesifikasi Kebutuhan Sistem}
\label{tab:acceptance}
\begin{tabularx}{\textwidth}{l p{3.2cm} c X}
\toprule
\textbf{Kode} & \textbf{Spesifikasi Kebutuhan} & \textbf{Hasil Uji} & \textbf{Bukti Verifikasi} \\
\midrule
SKF-01 & Multi-Provider API & TERPENUHI & Adaptor OpenAI, Claude, Gemini, Groq, dan Ollama terintegrasi. \\
SKF-02 & Mode REPL \& One-Shot & TERPENUHI & Eksekusi \texttt{crootagent} interaktif dan \texttt{crootagent "<task>"} berfungsi normal. \\
SKF-03 & Sesi Persisten & TERPENUHI & \texttt{test\_session\_serialization} lolos; berkas JSON tersimpan di disk. \\
SKF-04 & Operasi Berkas & TERPENUHI & \texttt{test\_write\_and\_read\_file} dan \texttt{test\_edit\_file} lolos 100\%. \\
SKF-05 & Eksekusi Shell & TERPENUHI & \texttt{test\_run\_command} berhasil mengeksekusi terminal dengan timeout. \\
SKF-06 & HITL Safety Confirmation & TERPENUHI & Prompt \texttt{[y/N]} muncul sebelum aksi destruktif; flag \texttt{-y} berfungsi. \\
SKNF-01 & RAM Idle < 20 MB & TERPENUHI & Pengukuran empiris mencatat konsumsi memori hanya 10,2 MB RAM. \\
SKNF-02 & Startup < 50 ms & TERPENUHI & Pengukuran waktu inisialisasi rata-rata 18 milidetik. \\
SKNF-03 & Distribusi Multi-OS & TERPENUHI & Matriks GitHub Actions menghasilkan biner Linux, macOS, dan Windows. \\
\bottomrule
\end{tabularx}
\end{table}
